% ============================================================================
%  EVNet Sentinel — Findings
%
%  The single source of prose for the research write-up. Every number here is
%  reproduced by a named script in the repository; the command is given beside
%  each table so a reviewer can regenerate it rather than trust it.
%
%  Build:  pdflatex EVNet_Sentinel_Findings.tex   (run twice for references)
% ============================================================================
\documentclass[11pt, a4paper]{article}

\usepackage[utf8]{inputenc}
\usepackage[T1]{fontenc}
\usepackage{lmodern}
\usepackage[margin=1in]{geometry}
\usepackage{xcolor}
\usepackage{booktabs}
\usepackage{tabularx}
\usepackage{hyperref}
\usepackage{microtype}

% ---- optional packages -----------------------------------------------------
% BasicTeX ships without enumitem and tcolorbox, and installing them needs
% sudo. The document degrades to plain lists and framed boxes when they are
% absent, so it compiles on a minimal install as well as on Overleaf.
\newif\ifhasenumitem
\IfFileExists{enumitem.sty}{\usepackage{enumitem}\hasenumitemtrue}{\hasenumitemfalse}
\newif\ifhastcb
\IfFileExists{tcolorbox.sty}{\usepackage{tcolorbox}\hastcbtrue}{\hastcbfalse}
\usepackage{caption}
\usepackage{amssymb}

% ---- shared palette with the other project documents ----------------------
\definecolor{primary}{HTML}{1E40AF}
\definecolor{secondary}{HTML}{0D9488}
\definecolor{darkgray}{HTML}{374151}
\definecolor{lightbg}{HTML}{F8FAFC}
\definecolor{alertred}{HTML}{B91C1C}

\hypersetup{
    colorlinks=true,
    linkcolor=primary,
    urlcolor=secondary,
    citecolor=primary,
    pdftitle={EVNet Sentinel --- Findings},
    pdfauthor={EVNet Sentinel Team}
}

\ifhasenumitem
  \newenvironment{tightitem}%
    {\begin{itemize}[leftmargin=1.4em, itemsep=0.3em]}{\end{itemize}}
  \newenvironment{tightenum}%
    {\begin{enumerate}[leftmargin=1.6em, itemsep=0.3em]}{\end{enumerate}}
\else
  \newenvironment{tightitem}{\begin{itemize}}{\end{itemize}}
  \newenvironment{tightenum}{\begin{enumerate}}{\end{enumerate}}
\fi

\ifhastcb
  \newtcolorbox{finding}[1]{
      colback=lightbg, colframe=primary, fonttitle=\bfseries,
      title={#1}, breakable, boxrule=0.8pt, arc=2pt
  }
  \newtcolorbox{caution}[1]{
      colback=white, colframe=alertred, fonttitle=\bfseries,
      title={#1}, breakable, boxrule=0.8pt, arc=2pt
  }
\else
  % Fallback: a ruled block with a coloured heading. Plainer, same information.
  \newenvironment{finding}[1]
    {\par\medskip\noindent\textcolor{primary}{\rule{\textwidth}{1pt}}\par\nobreak
     \smallskip\noindent\textbf{\textcolor{primary}{#1}}\par\nobreak\smallskip\itshape}
    {\par\smallskip\noindent\textcolor{primary}{\rule{\textwidth}{1pt}}\par\medskip}
  \newenvironment{caution}[1]
    {\par\medskip\noindent\textcolor{alertred}{\rule{\textwidth}{1pt}}\par\nobreak
     \smallskip\noindent\textbf{\textcolor{alertred}{#1}}\par\nobreak\smallskip\itshape}
    {\par\smallskip\noindent\textcolor{alertred}{\rule{\textwidth}{1pt}}\par\medskip}
\fi

\newcommand{\repro}[1]{\par\smallskip{\footnotesize\textcolor{darkgray}{%
    \textbf{Reproduce:} \texttt{#1}}}\par}
\newcommand{\col}[1]{\texttt{\small #1}}

\title{\textbf{EVNet Sentinel}\\[0.3em]
       \large Leakage-Corrected Intrusion Detection on CICEVSE2024}
\author{EVNet Sentinel Team}
\date{\today}

\begin{document}
\maketitle

% ===========================================================================
\begin{abstract}
\noindent
We reproduce the online intrusion-detection method of Makhmudov et al.\ (2025)
on the CICEVSE2024 electric-vehicle charging dataset and find that its reported
multiclass accuracy is an artifact of \emph{capture-session timestamp leakage}.
A decision tree given the single column \col{bidirectional\_first\_seen\_ms} and
nothing else recovers the 15-class label at $0.9976$, exceeding the published
$0.9840$. The reference preprocessing compounds the problem by discarding
\col{dst\_port} --- the feature that legitimately identifies a port scan --- while
retaining the timestamp that identifies it only by recording time. Removing both
fingerprints collapses $1{,}724{,}504$ reconnaissance flows to $6{,}821$.

With the leak removed, flow-level features separate volumetric denial-of-service
families near-perfectly (mean $F_1 = 1.000$, zero variance across seeds) but
cannot distinguish \texttt{nmap} reconnaissance modes (mean $F_1 = 0.183$), and
the two groups carry comparable support --- so the gap is a property of the data
rather than of model capacity. Separately, prequential accuracy in capture order
is inflated a further 44 points by label autocorrelation, and 42 of 43 detected
concept-drift events coincide with the seams where packet captures were
concatenated.
\end{abstract}

% ===========================================================================
\section{Dataset and Pipeline}

CICEVSE2024 is access-restricted and carries no public version string, so the
capture set is pinned by fingerprint rather than by version.

\begin{table}[h]
\centering
\caption{Dataset provenance.}
\begin{tabularx}{\textwidth}{lX}
\toprule
\textbf{Property} & \textbf{Value} \\
\midrule
Source            & CICEVSE2024, Canadian Institute for Cybersecurity (network traffic subset) \\
Capture files     & 59 \\
Raw flows         & 2{,}744{,}700 \\
Classes           & 15 (14 attack types + benign) \\
Aggregate SHA-256 & \texttt{\footnotesize be68ff18bbb86f217c7c0852c988f1fa4503b79bbda11c89960cf2700b63ea07} \\
\bottomrule
\end{tabularx}
\end{table}
\repro{python3 src/reproduction/dataset\_manifest.py --raw-dir data/raw}

\begin{table}[h]
\centering
\caption{Preprocessing stages, \texttt{--feature-set extended}.}
\begin{tabular}{lrr}
\toprule
\textbf{Stage} & \textbf{Rows} & \textbf{Columns} \\
\midrule
59 capture CSVs concatenated, labelled from filename      & 2{,}744{,}700 & 86 + label \\
After identifier, sparse, timestamp and \col{src\_port} drops & 2{,}744{,}700 & 65 \\
After deduplication (\texttt{--dedup global})              & \textbf{1{,}198{,}151} & 65 \\
Feature matrix (metadata and label separated)              & 1{,}198{,}151 & \textbf{60} \\
Train / validation / test (70/15/15)                       & 838{,}705 / 179{,}723 / 179{,}723 & 60 \\
\bottomrule
\end{tabular}
\end{table}
\repro{make data-process}

\medskip
\noindent
Labels are derived from the capture filename inside \col{preprocess.py}, porting
the reference implementation's 19-alias map. An earlier revision of this project
depended on a pre-labelled private archive and could not rebuild the dataset from
the public download; that step is now in version control.

% ===========================================================================
\section{Finding 1: The label is recoverable from a single timestamp column}

The reference preprocessing retains six absolute epoch-millisecond columns:
\col{bidirectional\_first\_seen\_ms}, \col{bidirectional\_last\_seen\_ms},
\col{src2dst\_first\_seen\_ms}, \col{src2dst\_last\_seen\_ms},
\col{dst2src\_first\_seen\_ms} and \col{dst2src\_last\_seen\_ms}.

Each attack in CICEVSE2024 was captured in its own packet capture at its own
wall-clock time, so these columns constitute a capture-session identifier rather
than a property of the traffic. None of the three drop stages in the reference
notebook removes them.

\begin{table}[h]
\centering
\caption{What a single feature recovers on its own. Decision tree,
         \texttt{max\_depth=25}, 15-class target, 3-fold cross-validation.}
\begin{tabular}{llr}
\toprule
\textbf{Feature} & \textbf{Nature} & \textbf{Accuracy} \\
\midrule
\col{bidirectional\_first\_seen\_ms} & capture fingerprint         & \textbf{0.9976} \\
\col{src\_port}                      & partial capture fingerprint & 0.4028 \\
\col{dst\_port}                      & legitimate signal           & 0.2917 \\
\midrule
\multicolumn{2}{l}{\emph{Published multiclass accuracy (Makhmudov et al.)}} & \emph{0.9840} \\
\bottomrule
\end{tabular}
\end{table}

\begin{finding}{A single-column decision stump exceeds the published result}
Replaying the reference preprocessing verbatim yields $1{,}277{,}520$ rows
against its reported ``$\sim$1.2M'', confirming the reconstruction is faithful.
On that feature set, a tree using \col{bidirectional\_first\_seen\_ms} alone
scores $0.9976$ --- higher than the $0.9840$ the paper reports for its full
model. In our own pre-correction pipeline, $39.8\%$ of the Random Forest's total
feature importance sat on those six columns.
\end{finding}
\repro{python3 src/reproduction/replay\_reference\_preprocessing.py --audit}

% ===========================================================================
\section{Finding 2: Reconnaissance rests on two columns}

Reconnaissance flows are distinguished almost entirely by \col{dst\_port} and the
capture timestamp. Holding the reference feature set fixed and toggling only
those two, then deduplicating as the reference does:

\begin{table}[h]
\centering
\caption{Ablation over the two fingerprints. Flood classes are unaffected
         throughout ($1{,}001{,}106 \rightarrow 970{,}079$, $96.9\%$ retained).}
\begin{tabular}{llrrr}
\toprule
\textbf{Timestamps} & \textbf{\col{dst\_port}} & \textbf{Total rows} & \textbf{Recon.\ flows} & \textbf{\% of max} \\
\midrule
kept    & kept    & 2{,}728{,}799 & 1{,}724{,}504 & 100.0\,\% \\
\textbf{kept} & \textbf{dropped} & \textbf{1{,}277{,}520} & \textbf{273{,}932} & \textbf{15.9\,\%} \\
dropped & kept    & 1{,}198{,}143 & 194{,}737 & 11.3\,\% \\
dropped & dropped & 977{,}862 & \textbf{6{,}821} & \textbf{0.4\,\%} \\
\bottomrule
\end{tabular}
\end{table}
\repro{python3 src/reproduction/ablate\_fingerprints.py}

\noindent
The emphasised row is the published preprocessing. Each column is independently
near-sufficient, so neither alone explains the collapse --- but the two are not
equivalent in kind:

\begin{tightitem}
  \item \col{dst\_port} is a \textbf{legitimate} discriminator. A port scan is
        \emph{defined} by sweeping destination ports; that is the behaviour, not
        an artifact.
  \item The capture timestamp is \textbf{not}. It identifies which recording the
        flow came from.
\end{tightitem}

\begin{finding}{The reference pipeline discards the real signal and keeps the artifact}
By dropping \col{dst\_port} while retaining the timestamps, the reference
preprocessing removes the feature that legitimately identifies reconnaissance and
keeps the one that identifies it only by accident of recording time. Its
$273{,}932$ reconnaissance flows are therefore separable by \emph{when they were
captured} rather than by what they did --- which is precisely why a stump on
\col{bidirectional\_first\_seen\_ms} reaches $0.9976$.
\end{finding}

\begin{caution}{Correction to an earlier revision}
An earlier version of this document attributed the reconnaissance collapse to the
timestamps alone. The ablation above shows that dropping \col{dst\_port} accounts
for the larger share. The finding stands; the mechanism as previously stated did
not.
\end{caution}

% ===========================================================================
\section{Finding 3: Stream order, and what the drift detector is detecting}

\subsection{Label autocorrelation}

Evaluated prequentially on an identical 25{,}000-instance slice with leak-free
features, varying only the order in which the stream arrives:

\begin{table}[h]
\centering
\caption{Stream ordering, controlled. \texttt{file-then-shuffle} uses the same
         instances and the same class distribution as \texttt{file}, differing
         only in order.}
\begin{tabular}{lrrr}
\toprule
\textbf{Order} & \textbf{Accuracy} & \textbf{Weighted $F_1$} & \textbf{Drift events} \\
\midrule
Capture order (\texttt{file})           & \textbf{0.9976} & 0.9975 & 3 \\
Shuffled, controlled (\texttt{file-then-shuffle}) & \textbf{0.5508} & 0.5318 & 1 \\
Shuffled, whole stream                  & 0.5029 & 0.4909 & 1 \\
\bottomrule
\end{tabular}
\end{table}

In capture order, consecutive flows share a label, so a prequential learner
scores near-perfectly by predicting ``the same as recently''. This is an
inflation mechanism entirely separate from the timestamp leak of Finding 1.

\subsection{Drift events are capture-file seams}

CICEVSE2024 is 59 packet captures concatenated end to end, and each join is an
abrupt distribution change by construction. Running ADWIN prequentially over the
full $1{,}921{,}182$-instance stream in capture order:

\begin{table}[h]
\centering
\caption{Alignment of detected drift events with capture-file boundaries,
         against a null model of 2{,}000 random placements.}
\begin{tabular}{lr}
\toprule
\textbf{Quantity} & \textbf{Value} \\
\midrule
Drift events detected                                & 43 \\
Events within 100 instances of a capture boundary    & \textbf{42} \\
Median distance to the nearest boundary              & 22 instances \\
Expected under random placement                      & $2.2 \pm 1.5$ \\
Empirical significance                               & $p < 0.0005$ \\
\bottomrule
\end{tabular}
\end{table}

\begin{finding}{The detector is finding the seams, not the traffic}
$42$ of $43$ drift events land within 100 instances of a point where one
recording was stitched to the next, against $2.2 \pm 1.5$ expected by chance.
These events are not evidence that charging-station traffic evolves; they mark
dataset assembly. This matters because the adaptivity of the online method is
justified in the original work by the presence of concept drift.
\end{finding}

% ===========================================================================
\section{Finding 4: \col{src\_port} is a residual fingerprint}

Ephemeral source ports are allocated near-sequentially by the operating system,
so within a single capture they occupy a contiguous band --- a lower-resolution
version of the timestamp fingerprint.

\begin{table}[h]
\centering
\caption{Controlled measurement: identical rows, identical split, column removed
         at fit time so deduplication cannot confound the comparison.}
\begin{tabular}{lrrr}
\toprule
\textbf{Model} & \textbf{With \col{src\_port}} & \textbf{Without} & \textbf{$\Delta$} \\
\midrule
Random Forest (depth 15, 100 trees) & 0.5894 & 0.5138 & $-0.0756$ \\
Decision Tree (depth 20)            & 0.6834 & \textbf{0.5135} & $\mathbf{-0.1699}$ \\
\bottomrule
\end{tabular}
\end{table}
\repro{python3 src/reproduction/srcport\_control.py}

\noindent
The deeper model loses more than twice as much, which is the signature of a
memorisable fingerprint rather than a genuine feature: additional depth buys
finer \col{src\_port} bands and nothing else. It also resolves an anomaly in the
first round of corrected results, where an unrestricted Decision Tree appeared to
outperform the ensemble. Once \col{src\_port} is removed the two converge
($0.5138$ versus $0.5135$), so that entire advantage was memorisation.

\col{dst\_port} is retained: unlike \col{src\_port} it is the behaviour that
defines a port scan.

% ===========================================================================
\section{Corrected Results}

\begin{caution}{Macro-$F_1$ is the headline, not accuracy}
With \texttt{ICMP\_Fragmentation} at 28 flows and \texttt{Benign} at 81, accuracy
tracks the volumetric floods and compresses genuinely different models into the
same number. The corrected Random Forest scores $0.8679$ accuracy against
$0.5582$ macro-$F_1$ on identical predictions.
\end{caution}

\begin{table}[h]
\centering
\caption{Multiclass results, mean $\pm$ standard deviation over seeds
         42 / 1337 / 2024. Each seed re-runs the train/test split \emph{and}
         model initialisation.}
\begin{tabular}{lrrrr}
\toprule
\textbf{Model} & \textbf{Macro $F_1$} & \textbf{Weighted $F_1$} & \textbf{Accuracy} & \textbf{Fit (s)} \\
\midrule
Random Forest       & \textbf{0.5582 $\pm$ 0.0080} & 0.8654 & 0.8679 & 8.6 \\
Decision Tree       & 0.5568 $\pm$ 0.0095 & 0.8677 & 0.8681 & 3.5 \\
Logistic Regression & 0.4355 $\pm$ 0.0172 & 0.8529 & 0.8626 & 71.7 \\
SVM (SGD, hinge)    & 0.4110 $\pm$ 0.0165 & 0.8489 & 0.8623 & 36.3 \\
\bottomrule
\end{tabular}
\end{table}
\repro{python3 src/evaluation/run\_experiments.py --raw-dir data/raw --feature-set extended}

\noindent
Random Forest and Decision Tree overlap at one standard deviation, independently
confirming Finding 4: their earlier apparent gap was \col{src\_port}
memorisation, not modelling.

\subsection{The central result, per class}

\begin{table}[h]
\centering
\caption{Per-class $F_1$, mean over three seeds. ``Support'' is the total flow
         count across all partitions after deduplication.}
\small
\begin{tabular}{llrrrrr}
\toprule
\textbf{Class} & \textbf{Group} & \textbf{Support} & \textbf{RF} & \textbf{DT} & \textbf{LogReg} & \textbf{SVM} \\
\midrule
PSHACK Flood        & Vol. & 195{,}952 & 1.000 & 1.000 & 1.000 & 1.000 \\
SynonymousIP Flood  & Vol. & 256{,}739 & 1.000 & 1.000 & 0.991 & 0.992 \\
TCP Flood           & Vol. & 256{,}297 & 1.000 & 1.000 & 1.000 & 1.000 \\
SYN Flood           & Vol. & 259{,}484 & 1.000 & 1.000 & 0.991 & 0.992 \\
UDP Flood           & Vol. & 32{,}072  & 0.999 & 0.999 & 1.000 & 1.000 \\
\midrule
TCP Port Scan       & Rec. & 35{,}123 & 0.265 & 0.202 & 0.278 & 0.100 \\
OS Fingerprinting   & Rec. & 27{,}702 & 0.223 & 0.227 & 0.039 & 0.022 \\
Vulnerability Scan  & Rec. & 39{,}002 & 0.204 & 0.230 & 0.215 & 0.201 \\
Service Version Det.& Rec. & 30{,}338 & 0.156 & 0.174 & 0.076 & 0.027 \\
Aggressive Scan     & Rec. & 27{,}815 & 0.145 & 0.179 & 0.081 & 0.041 \\
SYN Stealth Scan    & Rec. & 34{,}765 & 0.105 & 0.175 & 0.042 & 0.183 \\
\midrule
Benign              & Oth. & 81    & 0.986 & 0.941 & 0.225 & 0.192 \\
ICMP Flood          & Oth. & 32    & 0.660 & 0.638 & 0.216 & 0.178 \\
ICMP Fragmentation  & Oth. & 28    & 0.439 & 0.422 & 0.095 & 0.094 \\
Slowloris Scan      & Oth. & 2{,}721 & 0.191 & 0.166 & 0.283 & 0.145 \\
\bottomrule
\end{tabular}
\end{table}

\begin{finding}{The gap is a property of the data, not of model capacity}
Volumetric floods reach mean $F_1 = 1.000$ with \emph{zero variance across
seeds}; reconnaissance sits at $0.183$. The obvious objection --- that
reconnaissance simply has fewer examples --- does not hold: \texttt{UDP\_Flood}
has $32{,}072$ flows and scores $0.999$, while \texttt{TCP\_Port\_Scan} has
$35{,}123$ flows and scores $0.265$. Comparable support, opposite outcomes.

All six reconnaissance classes are \texttt{nmap} modes, and \texttt{-A} is
literally \texttt{-sV -O -sC --traceroute}, a superset of two of the others. At
flow level they emit near-identical single-probe flows and \emph{should} be
confusable. This establishes a realistic ceiling for flow-based EVCS intrusion
detection rather than a failure of the models.
\end{finding}

\noindent
The four classes in the ``Other'' group are excluded from both group means and
reported separately: three have fewer than 100 flows in the entire dataset, and
Slowloris is a low-rate attack whose flows resemble neither group.

\subsection{Effect of removing the leak}

\begin{table}[h]
\centering
\caption{Macro-$F_1$ before and after removing the six timestamp columns,
         like-for-like models.}
\begin{tabular}{lrrr}
\toprule
\textbf{Model} & \textbf{Leaked} & \textbf{Corrected} & \textbf{$\Delta$} \\
\midrule
Random Forest       & 0.9547 & 0.5894 & $\mathbf{-0.3653}$ \\
Logistic Regression & 0.4622 & 0.3308 & $-0.1314$ \\
SVM (SGD, hinge)    & 0.2829 & 0.3084 & $+0.0255$ \\
\bottomrule
\end{tabular}
\end{table}

\noindent
The linear models barely move, because a hyperplane cannot exploit a timestamp
axis the way a tree can. That --- not non-linearity --- is what the original
$0.9994$-versus-$0.3665$ gap was measuring.

% ===========================================================================
\section{Reproducibility}

\begin{table}[h]
\centering
\caption{Environment in which the reported numbers were produced.}
\begin{tabular}{ll}
\toprule
\textbf{Component} & \textbf{Value} \\
\midrule
CPU               & Apple M5 (10 logical cores) \\
Memory            & 16 GB \\
Python            & 3.13.12 \\
NumPy             & 2.4.4 \\
pandas            & 3.0.2 \\
scikit-learn      & 1.8.0 \\
River             & 0.26.1 \\
\bottomrule
\end{tabular}
\end{table}

\noindent
\texttt{requirements.txt} pins exact versions. River's API changed materially
around 0.21, and scikit-learn's \texttt{class\_weight} handling and
tree-splitting behaviour have shifted across minor releases, so unpinned versions
make these numbers unreproducible.

Model artifacts are deliberately not version-controlled: \col{src/api/main.py}
loads every \texttt{.pkl} in \texttt{saved\_models/} into its serving registry,
so a stale committed model is served silently as though it were current. Until
this was corrected, the models in that directory were trained on leaked features.

% ===========================================================================
\section{Limitations and Future Work}

\subsection{Acknowledged limitations}

\begin{tightitem}
  \item \textbf{The benign class is effectively absent.} The network-traffic
        subset contains 81 benign flows in total. A binary attack-versus-benign
        task is therefore answered by predicting ``attack'' unconditionally, and
        was removed from this work. Reported binary accuracies on this dataset,
        including the original $0.9913$, should be read with that in mind.
  \item \textbf{Grouped and temporal splits are implemented but not used for
        headline numbers.} \col{preprocess.py} supports holding out whole
        captures and splitting by time; the reported results use a random split
        for comparability with prior work.
  \item \textbf{Host-events and power-consumption domains are unexplored.}
        CICEVSE2024 also provides hardware performance counters and power traces,
        which may separate the reconnaissance classes that flow-level features
        cannot.
\end{tightitem}

\subsection{Planned: a sequence-model comparator}

A recurrent comparator (LSTM or GRU) is a natural next step, and would test
whether the reconnaissance ceiling established here is a limit of \emph{flow-level
features} or of \emph{non-sequential models}. Two design constraints follow
directly from the findings above and must be respected, or the comparator will
reproduce the artifacts this work set out to remove.

\begin{caution}{A sequence model over an ordered stream will re-introduce Finding 3}
Training an LSTM on flows in capture order hands it exactly the label
autocorrelation that inflates prequential accuracy by 44 points. Consecutive
flows share a label, so a model with memory can reach near-perfect accuracy by
learning ``predict what came recently'' --- and will appear to beat every static
model here. Any sequence model must be evaluated against the
\texttt{file-then-shuffle} control, and its reported gain measured relative to a
trivial ``repeat the last label'' baseline.
\end{caution}

\noindent
The defensible formulations are:

\begin{tightenum}
  \item \textbf{Intra-flow packet sequences.} Model the packet-level sequence
        \emph{within} a single flow rather than a sequence of flows. This is the
        formulation that could genuinely beat the ceiling, because it accesses
        structure that NFStream's aggregation discards --- the ordering of packet
        sizes and inter-arrival times that might separate one \texttt{nmap} mode
        from another. It requires re-extracting from the \texttt{.pcap} files
        rather than using the provided flow CSVs.
  \item \textbf{Grouped-split sequences of flows.} Keep flow sequences, but hold
        out whole captures (\texttt{--split grouped --dedup per-capture}) so the
        model cannot exploit position within a capture it has already seen.
        Report against the autocorrelation baseline.
\end{tightenum}

\noindent
Formulation 1 is the one worth pursuing: it asks a question the current results
cannot answer. Formulation 2 mainly measures how well a model memorises capture
structure, which this work has already shown to be the dominant confound in this
dataset.

% ===========================================================================
\section*{Artifact Index}
\addcontentsline{toc}{section}{Artifact Index}

\begin{table}[h]
\centering
\small
\begin{tabularx}{\textwidth}{lX}
\toprule
\textbf{Script} & \textbf{Produces} \\
\midrule
\col{src/data\_prep/preprocess.py}                       & The corrected dataset; feature-set, split, dedup and scaling options \\
\col{src/evaluation/run\_experiments.py}                 & Every headline number, multi-seed \\
\col{src/reproduction/replay\_reference\_preprocessing.py} & Faithful replay of the reference pipeline, with leak audit \\
\col{src/reproduction/ablate\_fingerprints.py}           & The Finding 2 ablation \\
\col{src/reproduction/srcport\_control.py}               & The Finding 4 controlled measurement \\
\col{src/reproduction/dataset\_manifest.py}              & Dataset fingerprint \\
\col{src/models/arfadwin/train\_arfadwin.py}             & Prequential ARF + ADWIN, drift log with capture provenance \\
\bottomrule
\end{tabularx}
\end{table}

\vfill
\noindent\rule{\textwidth}{0.4pt}\\
{\footnotesize\textcolor{darkgray}%
{Reference implementation:
\url{https://github.com/TATU-hacker/Intrusion_Detection_on_Electric_Vehicle_Charging_Systems}\\
Makhmudov, F., Kilichev, D., Giyosov, U., Akhmedov, F. (2025).
\emph{Online Machine Learning for Intrusion Detection in Electric Vehicle
Charging Systems.} Mathematics 13(5), 712. DOI: 10.3390/math13050712}}

\end{document}
